{\large\textbf{Evolution of Supermassive Black Holes Binary}}

\textbf{Introduction}

The concept of gravitational waves is one of the most impressive predictions of
Einstein's Theory of General Relativity. Gravitational waves are the space-time
ripples propagating with the speed of light similarly to electromagnetic waves.
Direct detection of gravitational waves is incredibly difficult, however, the
first signal was detected on September 14, 2015, by LIGO and VIRGO
collaborations.

Gravitational waves are emitted during the rapid motion of massive objects. The
most powerful source of gravitational waves is the merging of two Supermassive
Black Holes (SBH). Black holes predicted by Theory of General Relativity
represent extremely compact objects which might have very large masses. Other
specific properties of black holes will not be needed in the solution of this
problem.

In the generally accepted theory of galaxies' evolution, it is supposed that
there is SBH with the mass ranging from $10^5 - 10^9$ of Solar masses in the
galaxy's center. Galaxies are huge stellar systems containing $10^{10} - 10^{11}$
stars. During their evolution, two galaxies can collide and merge into one.
What happens to two SBHs initially located in their centers? The evolution of
the SBH binary system can be divided into \textit{three} main stages. At each
stage SBHs approach each other, although the underlying physical phenomena
differ. We will examine these phenomena separately in the first three parts of
the problem. In the fourth part, we will use the obtained relations to
calculate the total time of the SBH binary system evolution.

At the end of their evolution, two SBHs will eventually approach each other and
merge into a single black hole. The merging process lasts about an hour and is
accompanied by an intense burst of gravitational radiation. Future
observatories like LISA will be able to detect this gravitational radiation.
Still, the research on the SBH evolution is under way now, at the dawn of
gravitational-wave astronomy.

\textbf{General information}

1. Express all your numerical answers in parsecs (pc) for the distances and
giga-years (Gy) for the time intervals. We will use Solar mass ($M_S$) as a
reference mass. You might need these values:

\begin{quote}
$1\,\mathrm{pc} = 3.1 \times 10^{16}\,\mathrm{m}$,

$M_s = 2.0 \times 10^{30}\,\mathrm{kg}$,

$t_H = 13.7\,\mathrm{Gy}$ , age of the Universe,

$G = 6.67 \times 10^{-11}$ N $\times$ m$^2$/kg$^2$,

$c = 3.0 \times 10^{8}$ m/s.
\end{quote}

2. When you encounter the word ``\textbf{estimate}'', you are not demanded the
exact answer. It is sufficient to obtain a result that differs from the
accurate one, not more than by a factor of 10. On the contrary, when you
encounter the word ``\textbf{find}'', you are supposed to achieve the exact
answer. The word ``\textbf{calculate}'' asks you to bring the numerical answer.

3. Throughout the problem, assume every star in the galaxy to have the same
mass $m = M_S$.

4. Throughout the problem we will not take into account the effects of the
Theory of General Relativity except gravitational waves emission. All stars
and black holes are considered as point masses governed by Newton's
gravitation law.

\textbf{Part A. Dynamic Friction (1.6 points)}

In this part we shall study the simplified model of the galaxy. You can ignore
the velocities of the stars in the galaxy and assume the constant stellar
concentration $n$. The characteristic size of the galaxy is $R$. The stellar
concentration is small enough, so the stellar collisions are extemely rare and
negligible. Let us consider a SBH with the mass $M \gg m$ moving with the
velocity $v$ through the galaxy. Surprisingly, the SBH experiences nonzero
average force from the stars. This force slows the motion of the SBH and is
called the force of dynamical friction for this reason. This part is devoted to
the determination of this force.

\textbf{A.1}\quad Let us work in the SBH's reference frame and consider the
transit of one star with impact parameter $b$ (fig. 1). Assume that
\begin{equation}
b \gg b_1 = \frac{GM}{v^2}.
\tag{1}
\end{equation}
The angular deflection of the star $\alpha = kb_1/b$, where $k$ is some
coefficient. Find the value of $k$. If you cannot find $k$, assume $k = 1$
hereafter. \hfill 0.75pt

\begin{center}\includegraphics[width=0.62\linewidth]{figures/APhO_2017_Q2_fig1.png}\end{center}

Fig. 1: The deflection of a star by the SBH with mass $M$. The impact parameter
is $b$, the minimal distance between the star and the SBH is $r_m$.

\textbf{A.2}\quad Let Ox axis be directed along the SBH's velocity. Find the
momentum component $\Delta p_x$ transferred from the star to the SBH.
\hfill 0.25pt

\textbf{A.3}\quad Estimate the average force $F_{DF}$ acting on the SBH by
taking the average over impact parameter $b$ . Neglect the contribution of the
stars with impact parameters $b < b_1$ . Assume the SBH to reside in the
central part of a galaxy. Express$F_{DF}$ in terms of $M$, $v$, $R$, $G$ and
stellar density $\rho = mn$ . \hfill 0.4pt

\textbf{A.4}\quad As you obtained in the previous task, the expression for
$F_{DF}$ includes the factor $\log R/b_1$, which we will denote further as
$\log\Lambda$. Calculate the value of $\log\Lambda$ for
$M = 10^8\ M_S$, $R = 20\,\mathrm{kpc} = 20 \times 10^3\ pc$ and velocity
$v = 200\,\mathrm{km/s}$. \hfill 0.2pt

\textbf{Part B. Gravitational sligshot (3.0 points)}

In this part, we will consider the system of two SBHs with equal masses
$M \gg m$ located in the center of the galaxy. Let's call this system a
\textbf{SBH binary}. We will assume that there are no stars near the SBH
binary, each SBH has a circular orbit of radius $a$ in the gravitational field
of another SBH.

\textbf{B.1}\quad Find the orbital velocity $v_{bin}$ of each SBH. Find the
total energy $E$ of the SBH binary. Express it in terms of $a$, $G$ and $M$.
\hfill 0.25pt

There are a lot of stars at distances much larger than $a$ from the binary.
Stars travel along complex and diverse trajectories in the gravitational field
of the whole galaxy. The motion of the stars can be considered chaotic, like
the motion of the molecules in an ideal gas. Let us assume that stars'
velocities have equal magnitudes $\sigma \ll v_{bin}$ and their average mass
density is $\rho$. In this case dynamical friction is no longer affecting the
SBH binary and energy losses are caused by other phenomenon.

\textbf{B.2}\quad Let us solve a related problem. Let a star of mass $m$
transit by a point mass $M_2 \gg m$ being at rest. The minimal distance between
the star and the point during the transit is $r_m$. The velocity of the star at
large distance is $\sigma$. Find the exact value of impact parameter $b$.
\hfill 0.5pt

If a star approaches the SBH binary for a distance about $a$, it participates
in a complex 3-body interaction with the binary that almost always results in a
star being shot out with the velocity about $v_{bin}$ (the velocity of the star
at the large distance after interaction). We will call such a strong
interaction a collision of a star with the SBH binary. Acceleration and the
shot of the star after the collision is called ``\textbf{gravitational
slingshot}''.

\textbf{B.3}\quad Estimate the characteristic time $\Delta t$ between two
successive collisions of the SBH binary with stars. Take into account that
$\sigma \ll v_{bin}$. \hfill 1.0pt

\textbf{B.4}\quad Estimate the SBH binary energy loss rate $dE/dt$. Estimate
the radius variation rate $da/dt$. Express it in terms of $a$, $\rho$,
$\sigma$, $G$. \hfill 0.25pt

\textbf{B.5}\quad Let us denote the initial radius of the system as $a_1$.
Estimate the time $T_{SS}$ for the radius to decrease by a factor of 2 due to
``gravitational slingshot''. Calculate $T_{SS}$ for
$\sigma = 200\,\mathrm{km/s}$, $a_1 = 1\,\mathrm{pc}$,
$\rho = 10^4 M_S/\mathrm{pc}^3$. \hfill 1.0pt

\textbf{Part C. Emission of gravitational waves (1.0 points)}

In this part we shall study the SBH binary with equal masses which doesn't
interact with the stars. Even in this case the system loses the energy due to
gravitational waves emission. The energy loss rate due to gravitational waves
is
\begin{equation}
\frac{dE}{dt} = -\frac{1024}{5}\frac{G}{c^5}(\omega^3 I)^2,
\tag{2}
\end{equation}
where $\omega$ is angular velocity of the binary, and\quad $I = 2Ma^2$ is
quadrupole moment of the system.

\textbf{C.1}\quad Find the SBH binary radius variation rate $da/dt$ due to the
emission of gravitational waves. \hfill 0.2pt

When the orbit radius of the SBH binary $a$ becomes close to the gravitational
radius of the black hole:
\begin{equation}
r_g = \frac{2GM}{c^2},
\tag{3}
\end{equation}
two SBHs quickly merge.

\textbf{C.2}\quad Let us denote the initial radius of the system as
$a_2 \gg r_g$. Estimate the time $T_{GW}$ it takes for the SBH binary to shrink
to the radius about $r_g$ due to the emission of gravitational waves. Express
$T_{GW}$ as a function of $a_2$, $M$, $c$ and $G$. \hfill 0.7pt

\textbf{C.3}\quad Calculate the initial radius $a_H$ of the binary of SBHs
with equal masses $M = 10^8 M_S$ if it takes it the age of the Universe to
merge: $T_{GW} = t_H$. \hfill 0.1pt

\textbf{Part D. Full evolution (4.4 points)}

In this part we will use the results obtained above. Let us consider the
\textbf{real} astrophysical situation. Two galaxies having SBH of mass
$M = 10^8\ M_S$ in their centers merged into a new stellar system. Let the new
galaxy be spherically symmetrical with radius
$R = 20\,\mathrm{kpc} = 20 \times 10^3\,\mathrm{pc}$. Let us assume that
stellar density varies with radius $r$ to the galaxy center as
\begin{equation}
\rho(r) = \frac{\sigma^2}{4\pi G r^2},
\tag{4}
\end{equation}
where $\sigma = 200\,\mathrm{km/s}$.

\textbf{D.1}\quad Let the body move in circular orbit of radius $a < R$ in
gravitational field of the stars. Neglect the force of dynamical friction and
find the velocity $v$ of the body. \hfill 0.25pt

Immediately after the merging of galaxies two SBHs have arbitrary positions
inside the new galaxy and do not affect each other. Let's consider one SBH. We
assume it moves in a circular orbit of radius $a < R$ around the galaxy center
and slowly loses energy due to the dynamical friction.

\textbf{D.2}\quad Estimate the orbit radius variation rate $da/dt$. In part A
we ignored the velocities of the stars. Although stars are moving in the real
galaxy, not all of them have exactly the same speed $\sigma$ . Instead, the
speed of the stars is only of the order of $\sigma$ , and so is the relative
speed of SBH with respect to the stars, hence you can use the result obtained
in A.3 for estimation. You should use the density $\rho(r)$ from equation (4).
Assume $\log\Lambda$ to be a constant calculated in A.4. \hfill 0.75pt

After a certain time, two SBH will approach the center of the galaxy. Let two
SBH move in a circular orbit of radius $a$ around the center galaxy in the
gravitational field of the stars.

\textbf{D.3}\quad Estimate the critical radius $a_1$ at which gravitational
interaction between two SBHs is no longer negligible and calculate it. We will
say that at this moment two SBHs form a binary system (fig. 2). \hfill 0.3pt

\begin{center}\includegraphics[width=0.92\linewidth]{figures/APhO_2017_Q2_fig2.png}\end{center}

\begin{center}Fig. 2: The evolution of SBHs before and after the formation of the binary system\end{center}

\textbf{D.4}\quad Let us assume that after the merging of galaxies two SBHs
were at distances $a_0 = 2\,\mathrm{kpc} = 2 \times 10^3\,\mathrm{pc}$ from the
galaxy center. Calculate the time $T_1$ it takes for two SBH to form a binary
due to dynamical friction. \hfill 0.75pt

After forming the binary, two SBHs shoot away all the stars from the center of
the galaxy and stay there alone. Since this moment, dynamical friction becomes
ineffective and the binary starts to lose the energy because of the slingshot
effect. You can assume that the velocities of the stars around the binary are
$\sigma$ and the stellar density is $\rho_1 = \rho(a_1)$ from the equation (4).
Slingshot effect shrinks the radius of the system drastically and after some
time the system starts to lose energy mostly due to the radiation of
gravitational waves.

\textbf{D.5}\quad When the binary radius is less than some value $a < a_2$ the
energy loss is caused by gravitational waves emission. Estimate the $a_2$ value
and calculate it. \hfill 0.3pt

\textbf{D.6}\quad Estimate the time $T_2$ of the binary radius reduction from
$a_1$ to $a_2$ (the slingshot stage). Estimate the time $T_3$ of binary radius
reduction from $a_2$ to almost zero (the stage of gravitational waves
emission). \hfill 1.75pt

\textbf{D.7}\quad For the parameters given above, calculate the total time
$T_{ev}$ of two SBH evolution from galaxies merging to SBH merging.
\hfill 0.3pt

\textbf{Historical remark.} For a long time astrophysicists have been thinking
the SBH binary evolution stops at the slingshot stage, since the binary has
shot out all the stars with small impact parameters which might collide with
it. It appeared that two SBH would never merge. This fact was called
\textbf{the final parsec problem}.

Real galaxies have complicated asymmetric shapes. Few years ago it was found
that in galaxies of complex shapes the stars with small impact parameters
appear again and again. The SBH binary continues to lose energy, but slower
than our estimation gives. The final parsec problem was successfully solved.
